% Passes 5840--5847: collision-recovery addendum to the matrix/Fourier frontier.
\subsection*{Determinant-polar doily and the integral binary shadow}

The Pass~5832--5839 matrix packet survives mathematically, but its pass namespace
was contaminated by a parallel packet landed after reservation.  We therefore keep
its exact certificate intact and continue with a collision-recovery addendum.

On $V=M_2(\mathbb F_2)$ put $q(X)=\det X$.  Its polarization
\[
B(X,Y)=q(X+Y)+q(X)+q(Y)
\]
is nondegenerate alternating.  Hence the fifteen nonzero matrices are a concrete
$W(3,2)$ point set.  The rank split is
\[
15=9+6,
\]
where the nine rank-one matrices are the hyperbolic quadric
\[
Q^+(3,2)\cong \mathbb P^1(\mathbb F_2)\times\mathbb P^1(\mathbb F_2)
\]
(the $3\times3$ grid) and the six units are its complement.  Of the fifteen
isotropic lines, six lie wholly in the grid and nine meet it in exactly one point.
This is an incidence-level coordinatization of the standard two-qubit doily $9+6$
split; the nonzero ring elements here are not literally the projective-ring points
of the Saniga--Planat--Pracna construction.

The full linear determinant stabilizer has order
\[
72=2\cdot36,
\]
with the $36$-element left/right group preserving the two rulings and transpose
exchanging them.  Since $|Sp(4,2)|=720$, the grid orbit has size
\[
720/72=10,
\]
matching the ten grid hyperplanes of $W(3,2)$.  Affinely this gives
\[
|2^4:O^+(4,2)|=1152.
\]
The shared order with $W(F_4)$ is not an isomorphism: exact generation gives
\[
|Z(2^4:O^+(4,2))|=1,
\qquad
|Z(W(F_4))|=2,
\]
with central $-I$ in the latter.

The integral $W_9$ firewall also sharpens modulo two.  The point/heavy lattice
$A_3^3$ has mod-$2$ Gram rank $6$ and radical dimension $3$, while
$A_3\otimes A_3$ has rank $4$ and radical dimension $5$.  These equal the
$2$-torsion ranks of their discriminant groups
\[
(\mathbb Z/4)^3,
\qquad
(\mathbb Z/4)^4\times\mathbb Z/16,
\]
respectively.  Explicitly, the line radical is
\[
\{\mathbf 1\otimes a+b\otimes\mathbf 1:a,b\in\mathbb F_2^3\},
\qquad \dim=3+3-1=5,
\]
whereas the point/heavy radical has one fibre-constant parity vector in each
$A_3$ block.

Finally, over a finite field $\mathbb F_q$, rank-one Fourier labels in
$M_2(\mathbb F_q)$ number $(q-1)(q+1)^2$ and each has $q-1$
factorizations.  Thus the point-fibre Radon preimage has kernel
\[
\boxed{(q-2)(q-1)(q+1)^2},
\]
so $q=2$ is the unique multiplicity-free case.  The sharp binary $W_9$
identification must therefore not be extrapolated to larger fields without
quotienting this factorization kernel.

\paragraph{Boundary.}
All statements above concern exact finite incidence geometry, integral lattices,
finite Fourier/Radon transforms and finite reflection groups.  They do not identify
the selected $q=5$ W33 carrier with a physical two-qubit system, and the mod-$2$
radicals are not physical gauge fields or error syndromes without an independent
hardware/code map.
